% GENERATED FILE. Edit canonical lessons, then run npm run generate:books.
% canonical-source-hash: fnv1a64:9da44982cb07139e
% canonical-lessons: IT-C431-caparra, IT-C431-acconto, IT-C431-disdire, IT-C431-penale, IT-C431-anticipare

\chapter{Prenotare e disdire}
\label{ch:it-prenotare-e-disdire}
\hlchaptermodality{431}

\begin{chapteropening}
\textbf{By the end of this chapter:} \emph{I can handle a booking: the deposit, the advance, cancelling, the cancellation fee and bringing dates forward.}

Complete the last lesson of chapter 431: I can handle a booking: the deposit, the advance, cancelling, the cancellation fee and bringing dates forward.
\end{chapteropening}

\section[la caparra]{la caparra — the deposit}

\label{lesson:IT-C431-caparra}



\begin{quote}
Before the new one: say the Italian for the camper van, then the Italian for the mosquito net.
\end{quote}

\subsection*{You'll want to know: la caparra}
\textbf{la caparra} — \textquotedblleft{}the deposit\textquotedblright{}.

\begin{quote}
Per la camera chiedono una caparra del trenta per cento.
\end{quote}

\begin{quote}
\emph{For the room they ask for a thirty per cent deposit.}
\end{quote}

If you cancel, you may lose \emph{la caparra}.

\begin{cousinweb}
\emph{caparra} is built from \emph{capo} 'head, chief', inherited from Latin \emph{caput}, and \emph{arra} 'pledge', from Latin \emph{arrha}, which came into Latin from Greek (Lewis \& Short).
\end{cousinweb}

\begin{grammarlens}[title={What you've built}]
One more word for telling, explaining and describing at B1.
\end{grammarlens}

\subsection*{Your turn}
\begin{itemize}
\raggedright
  \item \emph{Say it:} \emph{la caparra}
  \item \emph{Say it:} \emph{la caparra}, once more
  \item \emph{Say it:} a sentence of your own with \emph{la caparra}
  \item \emph{Recall:} say the Italian for the camper van, then the Italian for the mosquito net, then say \emph{la caparra} again
\end{itemize}

\begin{tcolorbox}[breakable,colback=teal!4,colframe=teal!35!black,title={Before you move on}]
What does la caparra mean? (\textquotedblleft{}The deposit\textquotedblright{}.) Say it once more.
\end{tcolorbox}
\section[l'acconto]{l'acconto — the advance, the down payment}

\label{lesson:IT-C431-acconto}



\begin{quote}
Before the new one: say the Italian for the mosquito net, then the Italian for the deposit.
\end{quote}

\subsection*{You'll want to know: l'acconto}
\textbf{l'acconto} — \textquotedblleft{}the advance, the down payment\textquotedblright{}.

\begin{quote}
Ho pagato un acconto e il resto all'arrivo.
\end{quote}

\begin{quote}
\emph{I paid an advance and the rest on arrival.}
\end{quote}

\emph{In acconto} — on account.

\begin{cousinweb}
\emph{acconto} is built from \emph{a} 'towards' and \emph{conto} 'account', which is inherited from Late Latin \emph{computus}, from \emph{computāre} 'to reckon', built on \emph{putāre} (Lewis \& Short).
\end{cousinweb}

\begin{grammarlens}[title={What you've built}]
One more word for telling, explaining and describing at B1.
\end{grammarlens}

\subsection*{Your turn}
\begin{itemize}
\raggedright
  \item \emph{Say it:} \emph{l'acconto}
  \item \emph{Say it:} \emph{l'acconto}, once more
  \item \emph{Say it:} a sentence of your own with \emph{l'acconto}
  \item \emph{Recall:} say the Italian for the mosquito net, then the Italian for the deposit, then say \emph{l'acconto} again
\end{itemize}

\begin{tcolorbox}[breakable,colback=teal!4,colframe=teal!35!black,title={Before you move on}]
What does l'acconto mean? (\textquotedblleft{}The advance, the down payment\textquotedblright{}.) Say it once more.
\end{tcolorbox}
\section[disdire]{disdire — to cancel (a booking)}

\label{lesson:IT-C431-disdire}



\begin{quote}
Before the new one: say the Italian for the deposit, then the Italian for the advance.
\end{quote}

\subsection*{You'll want to know: disdire}
\textbf{disdire} — \textquotedblleft{}to cancel (a booking)\textquotedblright{}.

\begin{quote}
Abbiamo dovuto disdire la prenotazione.
\end{quote}

\begin{quote}
\emph{We had to cancel the booking.}
\end{quote}

Conjugated like \emph{dire}: \emph{ho disdetto}.

\begin{cousinweb}
\emph{disdire} is built from \emph{dis-} 'un-' and \emph{dire} 'to say', which is inherited from Latin \emph{dīcere} (Lewis \& Short): to unsay what was agreed.
\end{cousinweb}

\begin{grammarlens}[title={What you've built}]
One more word for telling, explaining and describing at B1.
\end{grammarlens}

\subsection*{Your turn}
\begin{itemize}
\raggedright
  \item \emph{Say it:} \emph{disdire}
  \item \emph{Say it:} \emph{disdire}, once more
  \item \emph{Say it:} a sentence of your own with \emph{disdire}
  \item \emph{Recall:} say the Italian for the deposit, then the Italian for the advance, then say \emph{disdire} again
\end{itemize}

\begin{tcolorbox}[breakable,colback=teal!4,colframe=teal!35!black,title={Before you move on}]
What does disdire mean? (\textquotedblleft{}To cancel (a booking)\textquotedblright{}.) Say it once more.
\end{tcolorbox}
\section[la penale]{la penale — the cancellation fee, the penalty}

\label{lesson:IT-C431-penale}



\begin{quote}
Before the new one: say the Italian for the advance, then the Italian for to cancel.
\end{quote}

\subsection*{You'll want to know: la penale}
\textbf{la penale} — \textquotedblleft{}the cancellation fee, the penalty\textquotedblright{}.

\begin{quote}
Se disdici tardi, paghi una penale.
\end{quote}

\begin{quote}
\emph{If you cancel late, you pay a fee.}
\end{quote}

As an adjective, \emph{penale} means criminal: \emph{il codice penale}.

\begin{cousinweb}
\emph{penale} is a learned borrowing from Latin \emph{poenālis} 'of punishment', from \emph{poena} 'penalty' (Lewis \& Short).
\end{cousinweb}

\begin{grammarlens}[title={What you've built}]
One more word for telling, explaining and describing at B1.
\end{grammarlens}

\subsection*{Your turn}
\begin{itemize}
\raggedright
  \item \emph{Say it:} \emph{la penale}
  \item \emph{Say it:} \emph{la penale}, once more
  \item \emph{Say it:} a sentence of your own with \emph{la penale}
  \item \emph{Recall:} say the Italian for the advance, then the Italian for to cancel, then say \emph{la penale} again
\end{itemize}

\begin{tcolorbox}[breakable,colback=teal!4,colframe=teal!35!black,title={Before you move on}]
What does la penale mean? (\textquotedblleft{}The cancellation fee, the penalty\textquotedblright{}.) Say it once more.
\end{tcolorbox}
\section[anticipare]{anticipare — to bring forward; to pay in advance}

\label{lesson:IT-C431-anticipare}



\begin{quote}
Before the new one: say the Italian for to cancel, then the Italian for the cancellation fee.
\end{quote}

\subsection*{You'll want to know: anticipare}
\textbf{anticipare} — \textquotedblleft{}to bring forward; to pay in advance\textquotedblright{}.

\begin{quote}
Possiamo anticipare l'arrivo a venerdì?
\end{quote}

\begin{quote}
\emph{Can we bring our arrival forward to Friday?}
\end{quote}

\emph{In anticipo} — early.

\begin{cousinweb}
\emph{anticipare} is a learned borrowing from Latin \emph{anticipāre} 'to take beforehand', built from \emph{ante} 'before' and \emph{capere} 'to take' (Lewis \& Short).
\end{cousinweb}

\begin{grammarlens}[title={What you've built}]
One more word for telling, explaining and describing at B1.
\end{grammarlens}

\subsection*{Your turn}
\begin{itemize}
\raggedright
  \item \emph{Say it:} \emph{anticipare}
  \item \emph{Say it:} \emph{anticipare}, once more
  \item \emph{Say it:} a sentence of your own with \emph{anticipare}
  \item \emph{Recall:} say the Italian for to cancel, then the Italian for the cancellation fee, then say \emph{anticipare} again
\end{itemize}

\begin{tcolorbox}[breakable,colback=teal!4,colframe=teal!35!black,title={Before you move on}]
What does anticipare mean? (\textquotedblleft{}To bring forward; to pay in advance\textquotedblright{}.) Say it once more.
\end{tcolorbox}
